% Full GPSS normal form and positive decomposition proofs.

\subsection{GPSS normal form, finite positive assembly and future translation}
\label{sec:annex-p09-gpss-special-normal-form}

\begin{annexitem}{annex:p9.gpss.kamp.normal.form.bridge}{thm:p09-gpss-kamp-normal-form-bridge}{Stage closure and the exact constructor translations}
\begin{proof}
The inclusion of stage zero in stage one is trivial. If
\(\operatorname{Bet}_{G,n-1}\subseteq\operatorname{Bet}_{G,n}\)
in every endpoint configuration, every constructor admitted at
stage \(n\) is admitted at stage \(n+1\): its point literals
are still allowed, and each regional child is allowed at the
larger stage. This proves increasing stages by induction,
simultaneously for bounded intervals, both rays and the whole
order. It does not add an unprinted constructor to the grammar.

For a signed atom reverse its sign. For a regional component
\(e\) or \(\neg e\), with \(e\) an admitted existential,
interchange those two forms. The resulting component is equivalent
to the logical negation of the old one, including when the old
form was negative. Thus De Morgan's laws and finite distribution
put any finite Boolean combination into disjunctive normal form
over the same component types. Empty conjunction and disjunction
are retained as true and false. A finite list of regional children
has a common stage, its maximum, because the stages are increasing.

For the bridge, equality in a strict linear order satisfies
\[
                 x=y\quad\Longleftrightarrow\quad
                          \neg(x<y)\land\neg(y<x).
\]
The implication from left to right is irreflexivity; the reverse
uses comparison of distinct points. Renaming \(q_i\) to
\(P_i\) and replacing equality by this formula therefore
preserves every first-order formula by induction through Boolean
operations and quantifiers. Bound variables are renamed apart
before substitution.

A GPSS point literal is a one-literal Kamp point conjunction.
Translate a positive GPSS regional constructor by retaining its
guard and existential quantifier and translating its children.
If the child list is empty, insert the distinguished starred
Kamp true formula at the witness. This is exactly a permitted
Kamp existential constructor. For a negative GPSS constructor,
use
\[
 \neg\exists z\left(g_{ab}(z)\land\bigwedge_j\gamma_j\right)
 \quad\Longleftrightarrow\quad
 \forall z\left(g_{ab}(z)\to\bigvee_j\neg\gamma_j\right).
 \tag{1}
\]
Complement each child by its sign reversal and translate it
recursively. For an empty disjunction insert the starred Kamp
false formula. These are the allowed universal constructors.
The left formal endpoint gives the Kamp before row; the right
formal endpoint gives the after row; two variable endpoints give
the between row. Their permissible child regions agree exactly.
Induction on the finite GPSS construction proves both membership
and semantic equivalence.

In the other direction, a starred Kamp point conjunction becomes
a conjunction of GPSS point literals. Its distinguished true
formula becomes the empty conjunction; its distinguished false
formula becomes the empty disjunction. For a Kamp existential,
first translate all children. Distribute their finite disjunctions
to obtain a finite disjunction of conjunctions of permitted GPSS
children. Existential quantification distributes over that finite
disjunction, and each resulting guarded existential is a GPSS
positive component at one greater than a common child stage.
An empty child conjunction remains an admitted zero-list
constructor. An empty outer disjunction stays false.

For a Kamp universal, negate its disjunctive matrix, convert
that finite Boolean combination of translated children to
disjunctive normal form by the sign operation above, and use
the reverse of (1). Its inner existential becomes a finite
disjunction of GPSS positive components. The final negation
is a conjunction of their negative forms, again an allowed
disjunctive-normal-form clause. Empty inner disjunction makes
the result true; an empty conjunction under a guarded
existential tests existence of a point in the region, not
automatic truth. This proves every universal case as well.

All translations recurse on supplied finite construction
histories and perform only finite Boolean distributions.
They preserve the endpoint configuration and use no new
unary letter. The starred Kamp constants are important here:
an arbitrary order formula is not being declared a GPSS
point literal. Complete prefixes translate by the equality
identity, so the proved Kamp normal form has a precise route
into the GPSS special form for nonempty parameter lists.
The direct GPSS construction below also proves the sentence
case. No later Kamp elementary formula with unrestricted
past dependence has entered this point-label grammar
\cite[ch.~III, pp.~43--56]{Kamp1968}
\cite[p.~170]{GabbayPnueliShelahStavi1980}.
\end{proof}
\end{annexitem}

\begin{annexitem}{annex:p9.gpss.special.normal.form}{thm:p09-gpss-special-normal-form}{The complete special-form induction}
\begin{proof}
First rename bound variables so that they are pairwise distinct
and disjoint from all free variables. Use the ordinary finite
prenex transformations of classical first-order logic: negation
interchanges the two quantifiers, and a quantifier can move past
a Boolean factor in which its variable is not free. Freshness
ensures those side conditions. A vacuous quantifier can be
removed because the domain is nonempty. We normalize the
quantifier-free matrix and then eliminate the original quantifiers
from the innermost one outward. At each stage the declared
parameter list includes all remaining free and prefix variables.

\paragraph{Weak order types and the atomic basis.}
For a finite nonempty parameter set, equality of assigned points
partitions the variables, and strict order uniquely orders the
blocks. Conversely the prefix for an ordered partition forces
exactly those equalities and strict comparisons. There are finitely
many such partitions. Their prefixes are pairwise incompatible
and exhaustive for every assignment, including assignments in
orders too small to realize some of the types.

Fix a type and replace each variable in a unary atom by its
block representative. Every equality or order atom in the
matrix now has a fixed true or false value under the prefix.
The remaining finite Boolean expression is in unary atoms
at the representatives. Put it in disjunctive normal form
and retain the prefix on each clause. Each local conjunct
is a signed unary literal and hence belongs to
\(\operatorname{At}_G\) at that representative. A true
Boolean branch has no local conjunct; a false branch contributes
no disjunct. Disjoin over all order types. This proves the
quantifier-free case, including formulas with no unary predicates.

\paragraph{Boolean closure at a fixed parameter list.}
Conjunction of two special normal forms distributes into
conjunctions of clauses. Different complete order types have
false conjunction; matching types simply combine their local
lists. To complement a normal form, group its clauses by every
complete type \(\pi\). Under that prefix, the original
formula is a finite disjunction of local conjunctions. Negate
this finite Boolean expression and distribute to disjunctive
normal form. The complement operation of the preceding lemma
keeps each point or regional component in its own permitted
class. Retain \(\pi\) on every resulting clause. A type
with no original clause has a false local disjunction, whose
complement is true and therefore has an empty local conjunction.
Disjoining over all types proves complement closure without
negated or incomplete prefixes. Disjunction is simply list union.
These arguments account for both empty Boolean folds.

\paragraph{Existential quantification and equality blocks.}
Suppose a normalized matrix has declared parameters
\(X\cup\{z\}\), with \(X\ne\varnothing\).
Distribute \(\exists z\) over its finite disjunction.
Consider one clause. If \(z\)'s equality block also
contains a surviving parameter, its value is forced by that
parameter. If \(z\) is not the representative, delete its
equality from the prefix; the local components already use
the surviving representative. If \(z\) is the representative,
choose a surviving variable of its block and replace every
free occurrence of \(z\) in the local components and prefix
by that variable. Their endpoint types remain valid: induction
on a regional constructor shows that renaming an endpoint
renames the same guard and child endpoint without changing
its class. Fresh bound variables prevent capture. Canonicalize
the surviving prefix in either case. The forward implication
deletes or substitutes an equality, and the reverse implication
assigns \(z\) the existing equal point. No existence of an
extra point is assumed in these two equality cases.

\paragraph{Existential quantification in the three strict regions.}
The remaining possibility is that \(z\) is a singleton
block. Remove it and call the residual prefix \(\rho\).
The possible positions and the complete list of local
components which can involve \(z\) are
\[
\begin{array}{c|c|l}
 \text{position}&\text{guard}&\text{dependent components}\\ \hline
 \text{before first }u&z<u&
   \operatorname{Bef}_G(z),\operatorname{At}_G(z),
                                      \operatorname{Bet}_G(z,u)\\
 \text{between adjacent }u,v&u<z<v&
   \operatorname{Bet}_G(u,z),\operatorname{At}_G(z),
                                      \operatorname{Bet}_G(z,v)\\
 \text{after last }u&u<z&
   \operatorname{Bet}_G(u,z),\operatorname{At}_G(z),
                                      \operatorname{Aft}_G(z)
\end{array}
\]
Write \(I\) for the conjunction of all independent local
components and \(D\) for all dependent ones. The old
prefix equals \(\rho\) conjoined with the indicated
guard, because the surviving representatives already have
their specified order. Hence the clause transforms exactly as
\[
 \exists z(\rho\land g\land I\land D)
       \quad\Longleftrightarrow\quad
                   \rho\land I\land\exists z(g\land D).
 \tag{2}
\]
The final factor is respectively a before, between or after
GPSS component. Choose a common finite stage for all regional
children, keep the point literals as individual children, and
use the next constructor stage. If \(D\) has no conjunct,
this is the admitted zero-list existential. It still requires
a point in the region. A first point, a last point, or an
empty strict gap therefore makes that factor false, as required.
All independent components remain legal at the residual
representatives and neighboring regions. Disjoin the results
for all clauses to finish the existential step.

\paragraph{Universal quantification and empty positions.}
For a universal use the exact identity
\[
                       \forall z\,N\longleftrightarrow
                                  \neg\exists z\,\neg N.
 \tag{3}
\]
The proved complement procedure normalizes \(\neg N\)
on the same full parameter list; the existential procedure
normalizes its projection; and complement on the surviving
list gives the final normal form. Thus every step in (3)
has an explicit finite construction in this grammar.

For a surviving prefix with \(s\) blocks, the existential
procedure examines the \(s\) equal-point positions and
the \(s+1\) strict regions, both rays included. A matrix
missing a disjunct in a strict position makes the complemented
matrix true in that position. Its existential then tests
whether the region is nonempty, and the final complement
requires it to be empty. A missing equal-point position
instead supplies an existing named counterexample and makes
the universal branch false. In particular for an independent
condition \(B\),
\[
 \forall z(g(z)\to B)
       \longleftrightarrow B\lor\neg\exists z\,g(z).
\]
The last component is a negative zero-list regional existential.
It gives exactly the vacuous case, rather than discarding an
empty ray or gap. This verifies the universal construction
in all positions, not merely when every region has a point.

\paragraph{No surviving parameter and the empty outputs.}
When the last original variable \(z\) is eliminated,
each matrix clause on \(\{z\}\) consists of its trivial
self-equality and a list from the before, point and after
classes at \(z\). Its existential projection is a
whole-order component: the guard has both endpoints formal
and all the listed children have the required types.
Finite disjunction distributes as before. Universal
projection uses (3); whole-order signed components also
admit sign reversal, so the complement is a finite
disjunction of their conjunctions. This proves the sentence
extension. A variable-free true or false Boolean abbreviation
is the empty conjunction or disjunction from the outset.

Empty output disjunctions have their declared semantic
parameter lists even though false itself has no free
occurrences. If actual occurrences of every nonempty
parameter list are required, choose one complete prefix
and conjoin both \(e\) and \(\neg e\), where
\(e=\exists z(z<y_1)\) is a before component at its
first representative. This is an always-false special
clause containing the prefix variables and introducing
no unary predicate. True uses the disjunction of all
prefixes with empty local lists.

Each step removes one original prenex quantifier and
creates finite regional construction histories. Newly
created regional quantifiers are not recursively treated
as original prefix quantifiers. Thus the process terminates.
Only original unary literals are repeated, complemented
or moved, and the empty folds require no additional
predicate. This proves effectivity, support and all
asserted cases of the special-form lemma in its own
GPSS grammar. The earlier constructor bridge explains
precisely how the established Kamp normal form provides
the corresponding alternate route
\cite[Lemma~1, p.~170]{GabbayPnueliShelahStavi1980}.
\end{proof}
\end{annexitem}

\begin{annexitem}{annex:p9.gpss.future.special.form}{thm:p09-gpss-future-special-form}{Removing the empty past of the induced tail}
\begin{proof}
Fix a model and a point assigned to \(x\). Form its
induced suborder \(T_{\geq x}\), restricting each unary
predicate to this tail. It is nonempty and its first point
is the assigned \(x\). Because every quantifier in the
given formula is guarded by its inclusive lower bound
\(y\geq x\), evaluation in the original model equals
evaluation of the formula with those guards removed in
the induced tail. This follows by induction: atoms and
Boolean connectives are unchanged, and the permitted
witnesses or universally tested points are exactly the
tail domain for each guarded quantifier.

Apply the proved normal-form construction to that
unguarded syntax on the one declared parameter \(x\).
The result has local conjuncts only at \(x\), before
\(x\), or after \(x\). In the induced tail, a
positive before component is false because its outer
existential asks for a point strictly before the first
point. A negative before component is true for the
same reason. Replace them by these constants and
perform the finite Boolean simplifications. The
remaining clauses use only
\(\operatorname{At}_G(x)\cup\operatorname{Aft}_G(x)\).

Each surviving after component has the same truth in
the full model and the induced tail. Its outer witness
is strictly after \(x\), and induction on its
regional construction confines every further quantified
point to that witness's subinterval or future region,
still inside the tail. Point literals are also unchanged.
Thus the resulting uniform syntax is equivalent to
the original future formula in the original model.
The inclusive bound permits a quantified point to
coincide with \(x\); these equal-point cases are
exactly why its atomic conjuncts remain. No information
about the deleted past is inserted into those atoms
\cite[future formulas and Lemma~1, p.~170]{GabbayPnueliShelahStavi1980}.
\end{proof}
\end{annexitem}

\begin{annexitem}{annex:p9.gpss.future.label.closure}{thm:p09-gpss-future-label-closure}{Future locality and finite Boolean labels}
\begin{proof}
Finite Boolean combinations of \(L(U)\)-formulas are
again in that language, and the first-order translation
commutes with each Boolean connective. The empty true
and false combinations have the displayed Boolean
encodings, whose values are independent of the chosen
inessential atom. This proves closure for every label
combination used later.

If two pointed models have the same ordered labelled
tail from the current point onward, their atomic values
agree there. Boolean constructors preserve agreement.
A strict until witness is later than the current
point, every point in its intervening interval is
also in that tail, and the operand evaluations at
all those points agree by induction. Thus the
until value agrees in both directions. This is
the complete \(L(U)\) induction and proves locality.

On \(\mathbb Z\), at current point zero, compare
the interpretation \(Q_0=\{-1\}\) with
\(Q_0=\varnothing\), keeping every other predicate
equal. Their tails from zero agree, so every
future label has the same value there. The formula
\(\exists v<u\,q_0(v)\), evaluated with \(u\mapsto0\),
has different values.
It therefore has no uniformly equivalent future
label. Being a one-variable first-order condition
alone does not satisfy the decomposition label
requirement
\cite[pp.~169--170]{GabbayPnueliShelahStavi1980}.
\end{proof}
\end{annexitem}

\begin{annexitem}{annex:p9.gpss.positive.decomposition}{thm:p09-gpss-positive-decomposition}{All finite merges and the positive regional constructor}
\begin{proof}
All equivalences in this proof are on the admissible
endpoint region. For finite endpoints that means
\(a<b\). Formal outer endpoints are treated by
omitting their impossible point quantification and
the corresponding vacuous comparisons. Point and
gap labels retain their whole-model interpretations.

\paragraph{Intersection of two decompositions.}
Suppose their interior chains have \(p\) and
\(q\) points. Enumerate every path from
\((0,0)\) to \((p,q)\) using steps
\((1,0),(0,1),(1,1)\). The first step type
places the next point of the first chain alone,
the second places the next point of the second
chain alone, and the third identifies their
next points. Each path determines one possible
weak merge of the two strictly ordered chains.
There are finitely many such paths, including
the empty one if both chains have no point.

For a merged point, conjoin the point labels
of all original chains which place a point
there. For each original chain which does
not, the point lies strictly in one of that
chain's gaps; conjoin the gap label at the
merged point as well. For a merged open
gap, conjoin the two original gap labels
covering it. These prescriptions are determined
solely by the finite merge pattern. The common
external endpoints receive no point labels,
since both original decompositions leave
them unlabelled. Every new label is a finite
conjunction of permitted future labels and
is therefore permitted by the label lemma.

If the upper endpoint is formal, carry both
finite lists of global cofinal conditions
into the output. Their conjunction is exactly
the list union; repetitions do not change
truth. The construction never changes those
conditions into a vacuous bounded-tail test.
Disjoin the decompositions obtained from
all merge paths.

To prove the forward implication, witnesses
for both original formulas have a definite
relative order, including every equality
between the chains. Their union realizes
one of the enumerated patterns. Original
point conditions hold at their own points;
an original gap condition holds also at
each newly inserted point in that gap.
On every smaller merged gap both covering
conditions hold. Thus the union witnesses
the constructed decomposition for that
pattern, with both cofinal lists intact.

Conversely, restrict any merged witness
chain to the positions belonging to each
original chain. The required point labels
are included among the output conjuncts.
Every point of an original open gap is
either a newly inserted merged point or
lies in one of its merged open subgaps.
In the first case its original gap label
was put at that point; in the second it
was put on the whole subgap. Therefore
the original universal gap condition
holds everywhere. This reasoning includes
empty subgaps and adjacent points. The
original cofinal conditions were retained,
so both original decompositions follow.
This proves intersection equivalence.

Finite disjunctions are handled by finite
distribution of conjunction over the two
lists, followed by this construction for
each pair. Repeating the binary construction
gives any finite nonempty conjunction.
For an empty conjunction use one gap
labelled true and no interior points.
It expresses the full admissible region,
including either ray and the whole-order
case. An empty disjunction denotes false
and contributes no branch to a product.

\paragraph{Concatenation at a quantified point.}
For a particular left and right decomposition,
their witnesses satisfy \(a<z<b\), with
all left internal points before \(z\)
and all right internal points after it.
Concatenate the left chain, the point
\(z\), and the right chain. Keep all
old point and gap labels and put
\(\theta\) at the new middle point.
The old decompositions do not label
their shared endpoint, so this adds
exactly the required condition there.
The new gaps are precisely the old
gaps in their order. If \(b=+\infty\),
the right decomposition's cofinal list
is carried to the output; the left
decomposition has an actual upper endpoint
and no such list.

Witnesses in either direction are
obtained by concatenating or splitting
this same finite chain at its designated
middle point. The outer strict guard
ensures that this is a genuine interior
point. No point is lost when one old
chain had no internal points, and a
gap with no available witness makes
both existential sides false. Distribute
over all pairs of input disjuncts to
obtain the assertion for finite lists.
This proof also applies when the left
endpoint is formal or both endpoints
are formal: their unbounded gap
universals remain exactly the same.

\paragraph{The first stage and the positive induction.}
A stage-one positive regional formula
has only a finite conjunction \(\theta\)
of point literals at its witness. This
is a future label. Use the decomposition
with one internal point labelled
\(\theta\), and both gaps labelled
true. Its witnesses are exactly the
original witnesses. The corresponding
negative regional formula is equivalent,
on the admissible interval, to the
decomposition with no internal points
and its sole gap labelled
\(\neg\theta\). This is exactly
the universal negation of the original
guarded existential. If the literal
list is empty, \(\theta\) is true:
the positive case requires an interior
point, and the negative case requires
the open interval to be empty. The
gap label false therefore does not
mean an always-false decomposition
on a discrete interval; it records
the genuine adjacency case.

For any higher positive regional
constructor whose child presentations
are supplied, combine all left-child
lists by the proved intersection
construction on \((a,z)\), and
all right-child lists on \((z,b)\).
An empty child list uses the true
one-gap unit on that region. Conjoin
the point literals at \(z\) into
the future label \(\theta\), using
true when there are none. Distribute
over the finitely many choices from
the left and right lists and apply
the concatenation construction to
each pair. If a list is empty as
a disjunction, no pair survives,
which correctly gives false.
These are all positive constructor
cases. A negative higher constructor
has not been silently treated as
positive: it is precisely where
the remaining complement proposition
is required.

Finally, each finite decomposition
contains its guard \(a<b\).
At \(a=b\), for example, the
negation of the true one-gap
decomposition is true, while every
decomposition disjunction is false.
Therefore an unguarded complement
statement cannot be valid on all
endpoint assignments. Restricting
the target to the admissible
interval preserves the source
normal-form use, whose prefix
already orders adjacent distinct
representatives
\cite[pp.~170--171]{GabbayPnueliShelahStavi1980}.

For \cref{def:p09-gpss-positive-merge-example}, two single
witnesses are either equal, strictly increasing or strictly
decreasing, giving exactly its three displayed point-label
lists. At endpoints zero and two in the three-point order,
there is only one interior point, so neither strict two-point
case is realizable. In the second example, the zero-point
decomposition's gap covers the other decomposition's middle
point. Its label \(q\) must therefore be conjoined at that
point as well as on the two new gaps. If \(p\) is true and
\(q\) false at the middle point, omitting that transferred
label would wrongly satisfy the merged formula, since both
new gap universals are vacuous. The required transfer prevents
this false result and agrees with the original conjunction.
\end{proof}
\end{annexitem}

\begin{annexitem}{annex:p9.gpss.suffix.translation}{thm:p09-gpss-suffix-translation}{Nested strict until and the complete cofinal condition}
\begin{proof}
First omit the cofinal list. At the
last internal point \(z_r\), the
remaining condition is that every
strictly later point satisfies the
last gap label. This is exactly
\(G\Phi_r\). Suppose the suffix
from internal point \(z_{i+1}\)
onward is expressed by \(H_{i+1}\).
From \(z_i\), choose the next
point strictly later, require
\(\Theta_{i+1}\land H_{i+1}\)
there, and require \(\Phi_i\)
at every intervening point. The
strict until clause gives exactly
\(\Phi_i\mathbin{\mathsf U_s}
(\Theta_{i+1}\land H_{i+1})\).
Induction backward through the
finite list proves the displayed
recursion. In either direction,
witnesses of each until give the
same strictly ordered points as
the decomposition. The label
translations keep their truth in
the original model, including
their own possible future queries
beyond subsequent witnesses.

If there is no internal point,
the decomposition has only its
future gap universal, giving
\(G\Phi_0\). At a last point
this universal is vacuous. If
an internal point is required,
the first strict until has no
witness at a last point, exactly
matching the decomposition.

For each global cofinal label,
we claim at any current
point \(x\in T\) that
\[
 \forall u\in T\ \exists v>u\ \psi(v)
       \quad\Longleftrightarrow\quad
                         F\Psi\land GF\Psi.
 \tag{4}
\]
The left side gives a later
\(\psi\)-point for \(u=x\),
and for each \(u>x\), proving
both right-hand conjuncts.
Conversely \(F\Psi\) gives
one \(v>x\) satisfying the
label. It is later than every
\(u\leq x\). For \(u>x\),
the \(GF\Psi\) conjunct
gives the required witness
after that \(u\). This
proves the global universal
over every domain point.
At a last point the first
conjunct is false, so (4)
does not inherit the vacuity
of \(GF\Psi\) alone.

Conjoin (4) for the finite
cofinal list with the already
proved formula \(H_0\).
An empty list contributes true.
All constructors and all label
formulas are in \(L(U)\) after
the stated Boolean/\(F,G\)
expansions. This proves the
entire positive suffix translation
at its actual endpoint scope
\cite[p.~171]{GabbayPnueliShelahStavi1980}.
\end{proof}
\end{annexitem}
